\section{Methodology}
\label{sec:methodology}

Our system is a pipeline of five stages. A feature builder joins the
data sources of Section~\ref{sec:data} into one table with a row per
(player, upcoming round). Three expected-points backends, a calibrated
heuristic, a structural independent-Poisson model, and a
gradient-boosted tree model, score that table independently; an
experimental Monte Carlo simulator provides a fourth, distributional
view. A per-position ensemble routes each position to the backend that
wins it on held-out data. A rule-encoding step converts predictions
into effective points, and a mixed-integer optimiser selects squads,
transfers, and lineups. Finally, a live layer supports the in-play
decisions the game allows after kickoff. This section specifies each
stage; formal statements of the optimisation programs and auxiliary
algorithms appear in Appendix~\ref{sec:b}.

Throughout, $i$ indexes a (player, round) row, $\pi(i) \in
\{\mathrm{GK}, \mathrm{DEF}, \mathrm{MID}, \mathrm{FWD}\}$ is the
player's position, $p_i$ is the price in millions, and $h_i \in \{0,
1\}$ indicates a home fixture. All backends share an availability
contract: a player whose status is not \emph{playing}, whose squad is
eliminated, or whom confirmed team news places on the bench is
predicted at zero.

\subsection{Per-player, per-round feature contract}
\label{sec:methodology-features}

The feature table contains one row per player per upcoming round in
which their squad plays. Its columns fall into five groups:

\begin{itemize}
\item \textbf{Player metadata:} identity, position, country, price,
  ownership fraction, availability status, an elimination flag, and
  the realised scoring history (cumulative points, last-round points,
  per-round points, and the game's own form value).
\item \textbf{Fixture context:} round, stage, kickoff time, home flag,
  opponent identity, and venue.
\item \textbf{Own-squad strength:} total and average squad price, and
  the mean price of the squad's top eleven players by price
  (\texttt{squad\_top\_n\_avg\_price}), which we use as a proxy for
  first-team strength.
\item \textbf{Opponent strength:} the same columns for the opposing
  squad, prefixed \texttt{opp\_}.
\item \textbf{Derived:} the top-eleven price gap
  (\texttt{strength\_diff}), the FIFA ranking gap
  (\texttt{rank\_diff}), the rolling country Elo of both sides and
  their difference (Section~\ref{sec:data-elo}), last-ten-match
  national-team form, and rest-day counts to the previous and next
  match.
\end{itemize}

The table is intentionally over-wide: each backend consumes only the
subset it needs, and adding a candidate feature never requires a
schema change downstream. Opponent strength is deliberately
first-class; fixture difficulty enters every backend rather than
being an afterthought.

\subsection{Backend 1: calibrated heuristic}
\label{sec:methodology-heuristic}

The first backend is a hand-calibrated formula with no training step.
Each row receives

\begin{equation}
\hat{y}_i \;=\; c_{\pi(i)}\, p_i
  \left(1 + \alpha \tanh z_i\right)
  \left(1 + \beta\, h_i\right)
  \;+\; \gamma \max\!\left(0,\; p_i - 9.0\right),
\label{eq:heuristic}
\end{equation}

where $c_{\pi}$ is a per-position points-per-price coefficient
($c_{\mathrm{GK}} = 0.50$, $c_{\mathrm{DEF}} = 0.55$,
$c_{\mathrm{MID}} = 0.60$, $c_{\mathrm{FWD}} = 0.65$), $\alpha = 0.40$
saturates the matchup adjustment at $\pm 40$\,\%, $\beta = 0.05$ is a
home boost, and $\gamma$ is a premium-tier boost above the 9.0M price
threshold, zero by default. The matchup signal $z_i$ blends a
price-based strength gap with the best available national-team
strength signal:

\begin{equation}
z_i =
\begin{cases}
0.35\, z^{\mathrm{price}}_i + 0.65\, z^{\mathrm{elo}}_i
  & \text{if country Elo is available,}\\[2pt]
0.35\, z^{\mathrm{price}}_i + 0.65\, z^{\mathrm{rank}}_i
  & \text{else if a FIFA ranking is available,}\\[2pt]
z^{\mathrm{price}}_i & \text{otherwise,}
\end{cases}
\label{eq:heuristic-z}
\end{equation}

with $z^{\mathrm{price}}_i = \Delta^{\mathrm{price}}_i / 2$,
$z^{\mathrm{elo}}_i = \Delta^{\mathrm{elo}}_i / 400$, and
$z^{\mathrm{rank}}_i = \Delta^{\mathrm{rank}}_i / 250$, where
$\Delta^{\mathrm{price}}$ is the top-eleven price gap,
$\Delta^{\mathrm{elo}}$ the country Elo gap, and
$\Delta^{\mathrm{rank}}$ the FIFA ranking-points gap. The divisors
normalise each signal to a comparable scale; 400 Elo points is the
rating system's natural one-tier distance \citep{elo1978rating}. The
Elo signal takes priority because it is derived live from realised
international results \citep{hvattum2010using}; the static ranking is
the fallback, and club-price alone carries the signal when neither is
present (as in club-league training data).

The heuristic's strengths are transparency and robustness: every
constant is explainable, no training data is required, and there is
no distribution-shift risk between club and tournament data. Its
weakness is that it cannot learn. If a real pattern exists that the
formula does not encode, the heuristic will never find it, and its
linear-in-price base term flattens the true, non-linear marginal
return of premium players. An attempt to graft a fixed-weight
realised-form blend onto this formula lost to the unmodified version
in every backtested round (Section~\ref{sec:negative-results}), which
informed the design of the learned form feature in
Section~\ref{sec:methodology-gbm}.

\subsection{Backend 2: structural Poisson model}
\label{sec:methodology-poisson}

The second backend follows the independent-Poisson tradition of
football score modelling \citep{maher1982modelling,
dixon1997modelling}: estimate each team's expected goals, distribute
them across positions, and map the result to fantasy points through
the official scoring rules. Like the heuristic it requires no
training data.

\paragraph{Step 1: team expected goals.}
A matchup score combines the available strength signals,

\begin{equation}
m_i = s_i + \frac{\Delta^{\mathrm{price}}_i}{6},
\qquad
s_i =
\begin{cases}
\Delta^{\mathrm{elo}}_i / 400 & \text{if Elo is available,}\\[2pt]
\Delta^{\mathrm{rank}}_i / 800 & \text{else if a ranking is
  available,}\\[2pt]
0 & \text{otherwise,}
\end{cases}
\label{eq:poisson-matchup}
\end{equation}

and maps to expected goals for the player's team and its opponent
through a clipped exponential around a base rate of 1.30 goals per
team per match, with a 10\,\% boost for the home side:

\begin{align}
\lambda^{\mathrm{own}}_i &= 1.30\,
  e^{\operatorname{clip}(m_i,\, [-1.2,\, 1.2])}
  \left(1 + 0.10\, h_i\right),
\label{eq:poisson-own}\\
\lambda^{\mathrm{opp}}_i &= 1.30\,
  e^{-\operatorname{clip}(m_i,\, [-1.2,\, 1.2])}
  \left(1 + 0.10\, (1 - h_i)\right).
\label{eq:poisson-opp}
\end{align}

Clipping the exponent keeps predictions sane for extreme mismatches.

\paragraph{Step 2: positional distribution.}
Team expected goals are distributed across positions with fixed
shares (Table~\ref{tab:poisson-constants}): a player's expected goals
are $\lambda^{\mathrm{own}}_i\, s^{G}_{\pi(i)}$ and, taking roughly
70\,\% of goals to carry an assist, expected assists are $0.7\,
\lambda^{\mathrm{own}}_i\, s^{A}_{\pi(i)}$.

\begin{table}[t]
\centering
\caption{Poisson backend constants per position: goal share $s^{G}$,
assist share $s^{A}$, points per goal $G$, and clean-sheet points
$C$. Point values follow the official game rules.}
\label{tab:poisson-constants}
\begin{tabular}{lcccc}
\toprule
Position & $s^{G}$ & $s^{A}$ & $G$ & $C$ \\
\midrule
GK  & 0.00 & 0.01 & 9 & 5 \\
DEF & 0.10 & 0.18 & 7 & 5 \\
MID & 0.30 & 0.50 & 6 & 1 \\
FWD & 0.50 & 0.31 & 5 & 0 \\
\bottomrule
\end{tabular}
\end{table}

\paragraph{Step 3: defensive terms.}
Under the Poisson assumption the clean-sheet probability is exact,
$P(\text{clean sheet}) = e^{-\lambda^{\mathrm{opp}}_i}$, and the
expected goals conceded beyond the first, which the rules penalise,
has the closed form

\begin{equation}
\xi_i = \mathbb{E}\bigl[\max(0,\, K - 1)\bigr]
      = \lambda^{\mathrm{opp}}_i - 1 + e^{-\lambda^{\mathrm{opp}}_i},
\qquad K \sim \mathrm{Poisson}\bigl(\lambda^{\mathrm{opp}}_i\bigr).
\label{eq:poisson-xi}
\end{equation}

\paragraph{Step 4: expected fantasy points.}
The components sum, with a 2-point appearance baseline, 3 points per
assist, a goalkeeper save bonus scaling at 0.50 points per unit of
opponent expected goals (an empirically calibrated constant; the
value derived from first principles measured worse, see
Section~\ref{sec:negative-results}), and the conceded-goals penalty
applied fully to goalkeepers and at rate 0.5 to defenders:

\begin{align}
\hat{y}_i = \max\Bigl(0,\;
  & 2
  + \lambda^{\mathrm{own}}_i\, s^{G}_{\pi(i)}\, G_{\pi(i)}
  + 3 \cdot 0.7\, \lambda^{\mathrm{own}}_i\, s^{A}_{\pi(i)}
  + e^{-\lambda^{\mathrm{opp}}_i}\, C_{\pi(i)} \nonumber\\
  & + \mathbf{1}\bigl[\pi(i) = \mathrm{GK}\bigr]
      \bigl(0.50\, \lambda^{\mathrm{opp}}_i - \xi_i\bigr)
    - \mathbf{1}\bigl[\pi(i) = \mathrm{DEF}\bigr]\, 0.5\, \xi_i
  \Bigr).
\label{eq:poisson-total}
\end{align}

Every term traces to a single football statement, and the clean-sheet
term is exact under the model. The weaknesses are the independence
assumption itself, which ignores score correlation between the two
teams \citep{karlis2003analysis}, and the population-average shares:
the model cannot tell a star striker from his backup if they share a
price and a matchup.

\subsection{Backend 3: gradient-boosted trees}
\label{sec:methodology-gbm}

The third backend is a LightGBM model \citep{ke2017lightgbm} per
position, reflecting the continued strength of tree ensembles on
tabular data at this scale \citep{shwartz2022tabular,
borisov2022deeptabular}. Each position trains four heads on the same
eight features:

\begin{equation}
\hat{y}_i = f^{\mathrm{LGBM}}_{\pi(i)}\!\left(
  p_i,\; h_i,\; \Delta^{\mathrm{price}}_i,\;
  \bar{p}^{\,11}_i,\; \bar{p}^{\,11,\mathrm{opp}}_i,\;
  \phi^{\mathrm{lag}}_i,\; \xi^{+}_i,\; \xi^{-}_i
\right),
\label{eq:gbm}
\end{equation}

where $\bar{p}^{\,11}$ and $\bar{p}^{\,11,\mathrm{opp}}$ are the
own-squad and opponent top-eleven average prices,
$\phi^{\mathrm{lag}}_i$ is a lagged form feature defined below, and
$\xi^{+}_i, \xi^{-}_i$ are the team's trailing real expected goals
for and against, added in the fourth model generation after passing
the walk-forward gate of Section~\ref{sec:validation-walkforward}
(they exist only for tournament rows; league training rows carry
missing values, which the trees handle natively). The
\emph{mean} head is an L2 regression on realised fantasy points; the
remaining heads are quantile regressions minimising the pinball loss

\begin{equation}
L_{\tau}(y, q) = \max\bigl(\tau\, (y - q),\; (\tau - 1)(y - q)\bigr),
\qquad \tau \in \{0.10,\; 0.50,\; 0.90\},
\label{eq:pinball}
\end{equation}

so the backend emits a point estimate together with a per-player
10th--90th percentile band; the upper quantile later feeds the
ceiling-aware captain choice (Appendix~\ref{sec:b}). Training data
are three seasons of English Premier League fantasy scoring
(Section~\ref{sec:data-fpl}); hyperparameters were selected by a
held-out sweep (15 leaves, learning rate 0.05, 200 boosting rounds,
minimum 30 samples per leaf, feature and bagging fractions 0.9 with
bagging every 5 iterations). Fixed seeds and LightGBM's deterministic
mode make training bit-reproducible; without them, run-to-run bagging
randomness perturbs held-out error enough to mask small
feature-engineering comparisons. Negative predictions are clipped to
zero.

\paragraph{Lagged form.}
The first two generations of this model were form-blind: all five of
their features were price or team-strength signals, so a player's
own recent output never entered the model, and in-form premium
players from outside the training league were systematically
under-predicted at the tournament. The current generation adds
$\phi^{\mathrm{lag}}_i$, the trailing mean of the player's realised
fantasy points over the previous three matches, shifted by one match
so the current row's label never enters its own feature. One
definition is shared by all three places the model reads data: club
training, tournament label extraction, and tournament inference. A
player's first match leaves the feature missing, which LightGBM
handles natively.

\paragraph{In-tournament retraining.}
The form feature alone is insufficient: a model fitted only on club
data learns an output scale that does not fit the tournament, and no
single feature corrects a scale mismatch. The second half of the fix
is retraining on the union of the club corpus and all completed
World Cup rounds, which recalibrates the output scale; the production
scheduler retrains before each scoring run, so realised tournament
labels can no longer accumulate unused. In the leak-free walk-forward
protocol of Section~\ref{sec:validation}, where the model for
held-out round $k$ trains only on club data plus tournament rounds
strictly before $k$, pooled RMSE improves from
\WalkforwardPooledA{} (form-blind, club-only) to
\WalkforwardPooledC{} (form feature plus tournament retraining).

The GBM's residual weakness is extrapolation: inputs unlike anything
in training leave a tree ensemble with nothing to draw on, which is
exactly the failure the retraining loop exists to bound.

\subsection{Monte Carlo backend (experimental)}
\label{sec:methodology-montecarlo}

A fourth backend simulates each match rather than computing a point
estimate. Three properties motivate it: it produces a full per-player
points distribution, which matters more in single-match knockout
rounds than over multi-match group horizons; it anchors on realised
tournament output through a per-player form multiplier, which the
heuristic and the pre-retraining GBM lacked; and it recalibrates
immediately as each round completes, with no retraining step.

For each row the simulator draws $N = 1000$ match realisations from
the Poisson backend's team-level estimates
(Equations~\ref{eq:poisson-own}--\ref{eq:poisson-opp}). In
realisation $s$,

\begin{align}
\ln \varepsilon^{\mathrm{own}}_s,\ \ln \varepsilon^{\mathrm{opp}}_s
  &\sim \mathcal{N}\bigl(0,\; 0.15^2\bigr),
\nonumber\\
G^{\mathrm{team}}_s \sim \mathrm{Poisson}\bigl(
  \lambda^{\mathrm{own}} \varepsilon^{\mathrm{own}}_s\bigr),
  &\qquad
G^{\mathrm{opp}}_s \sim \mathrm{Poisson}\bigl(
  \lambda^{\mathrm{opp}} \varepsilon^{\mathrm{opp}}_s\bigr),
\label{eq:mc-teams}\\
g_{i,s} \sim \mathrm{Poisson}\bigl(
  G^{\mathrm{team}}_s\, s^{G}_{\pi(i)}\, \phi_i\bigr),
  &\qquad
a_{i,s} \sim \mathrm{Poisson}\bigl(
  0.7\, G^{\mathrm{team}}_s\, s^{A}_{\pi(i)}\, \phi_i\bigr),
\label{eq:mc-players}
\end{align}

where the multiplicative log-normal noise captures form and news
uncertainty in team strength, and the Poisson means in
Equation~\ref{eq:mc-players} are clipped to guard the tails. The
form multiplier

\begin{equation}
\phi_i = \operatorname{clip}\!\left(
  \frac{\text{realised points per match of player } i}
       {\text{positional mean of the same quantity}},\;
  [0.3,\; 3.0]\right)
\label{eq:mc-form}
\end{equation}

is the realised-data anchor: a player producing at twice his
positional average receives twice the goal and assist share.
Goalkeeper saves are sampled as $\mathrm{Binomial}(\text{shots},
0.70)$ with shots drawn from $\mathrm{Poisson}(2.0\,
\lambda^{\mathrm{opp}} \varepsilon^{\mathrm{opp}}_s)$, worth one
point per three saves. Each realisation is scored with the official
rules: 2 appearance points, per-position goal points, 3 per assist,
clean-sheet points when $G^{\mathrm{opp}}_s = 0$, and a defender
penalty of $0.5 \max(0,\, G^{\mathrm{opp}}_s - 1)$. The output is
the mean and the 10th, 50th, and 90th percentiles of the simulated
points, from a fixed seed for reproducibility.

The simulator deliberately omits minutes modelling (every row assumes
a 60-minute appearance) and in-match game-state dynamics. Its status
is experimental: the mean has not passed the held-out validation the
other backends were subjected to, so it is excluded from the
production run, but its distributional output informed the captaincy
work and it is retained as a parallel backend for comparison.

\subsection{Per-position ensemble}
\label{sec:methodology-ensemble}

Held-out validation (Section~\ref{sec:validation}) shows that no
single backend wins every position: the Poisson model wins
goalkeepers, the heuristic wins defenders, and the GBM wins
midfielders and forwards. The differences are small but stable
across runs. The production backend acts on this result by routing:

\begin{equation}
\hat{y}(i) = \hat{y}_{r(\pi(i))}(i),
\qquad
r = \bigl\{\mathrm{GK} \mapsto \text{Poisson},\;
  \mathrm{DEF} \mapsto \text{heuristic},\;
  \mathrm{MID},\, \mathrm{FWD} \mapsto \text{GBM}\bigr\}.
\label{eq:ensemble}
\end{equation}

This is not a new model but a refusal to keep using each model where
our own numbers say it is worst. The default routing encodes the
held-out winners and can be re-derived mechanically from any
validation report by taking the lowest per-position RMSE. The GBM's
quantile columns are carried through for the positions it serves;
positions routed to the structural backends are point estimates. Each
prediction row records which backend produced it, for auditability.

Two further reasons justify maintaining parallel backends rather than
a single model. First, the optimiser consumes exactly one prediction
column, and blending predictions shifts its optimum in non-obvious
ways; per-position routing is honest about which model we trust
where. Second, the backends cross-check each other: when they agree
we act with confidence, and when they diverge sharply, as they did
for premium forwards outside the GBM's training distribution, the
disagreement itself is diagnostic information. That disagreement is
what exposed the form-blindness defect fixed in
Section~\ref{sec:methodology-gbm}.

\subsection{Scouting bonus and effective points}
\label{sec:methodology-scouting}

The game awards a $+2$ scouting bonus when a player scores more than
4 points while owned by fewer than 5\,\% of teams. We encode the rule
as a deterministic post-prediction step applied exactly once,
downstream of every backend, using the prediction as a proxy for the
realised outcome:

\begin{equation}
b_i = 2 \cdot \mathbf{1}\bigl[\hat{y}_i > 4\bigr]
        \cdot \mathbf{1}\bigl[\omega_i < 0.05\bigr],
\qquad
\tilde{y}_i = \hat{y}_i + b_i,
\label{eq:scouting}
\end{equation}

where $\omega_i$ is the ownership fraction and $\tilde{y}_i$ the
effective points the optimiser consumes. The deterministic proxy
over-credits players predicted just above the threshold (who may
realise below it) and under-credits players just below it (who may
still trigger through realised variance), but the induced ordering
is adequate for squad selection. Applying the bonus once, centrally,
also prevents double counting: an early version baked it into one
backend's predictions and the downstream step re-applied it,
inflating every low-ownership backup goalkeeper. A
participation-based availability factor then multiplies effective
points before optimisation, discounting players whose tournament
history signals rotation risk (Appendix~\ref{sec:b}).

\subsection{Squad optimisation}
\label{sec:methodology-milp}

Squad decisions are mixed-integer linear programs, formulated in PuLP
\citep{pulp} and solved with the bundled CBC solver \citep{cbc}.
Three entry points share one constraint core; the full programs are
stated in Appendix~\ref{sec:b}.

\begin{itemize}
\item \textbf{Squad selection} chooses a 15-player squad maximising
  the sum of effective points over the stage horizon, subject to
  position counts (2 GK, 5 DEF, 5 MID, 3 FWD), the stage budget, the
  per-country cap, and exclusion of eliminated squads. The backup
  goalkeeper is discounted to his automatic-substitution value via a
  designated-scorer variable (Appendix~\ref{sec:b}): only one of the
  two mandatory goalkeepers can score in a round, so the optimum is
  one strong starter plus the cheapest legal backup.
\item \textbf{Transfer planning} adds the current squad as state and
  a continuous slack variable counting transfers beyond the free
  quota, each penalised at the official 3 points, so the solver
  trades squad improvement against transfer hits in one objective.
\item \textbf{Lineup selection} picks the starting eleven from the
  held squad for a single round: exactly one goalkeeper, one of the
  seven legal formations selected by binary indicator, and outfield
  counts tied to the chosen formation. The bench is ordered by
  predicted points within position. The solver's captain is a
  mean-argmax placeholder; production overrides it with a
  ceiling-aware selector that interpolates between the mean and the
  q90 head of Section~\ref{sec:methodology-gbm} according to the
  manager's standings percentile (Appendix~\ref{sec:b}).
\end{itemize}

Stage configurations encode the game's evolving rules: the budget
(\$100M in the group stage, \$105M from the round of 32), the
per-country cap (3 in the group stage, then 4, 5, 6, and 8 through
the final), and the free-transfer quota per stage. The two unlimited
transfer windows (pre-tournament and before the round of 32) are
modelled by short-circuiting transfer planning to a fresh squad
solve.

\subsection{Live decision tools}
\label{sec:methodology-live}

The game permits two in-play interventions, and a live module
supports both with expected-value calculations over the remaining
fixtures of a matchday.

\paragraph{Captain switch.}
Given the running scores of all fixtures in a matchday and the
deadline of the manual substitution option, the tool computes the
expected gain of moving the captaincy from the current pick to any
member of the starting eleven who has not yet played, comparing the
captain's realised points against the doubled prediction of each
not-yet-played alternative.

\paragraph{Substitution advisor.}
When a starter finishes a match with zero minutes, the game's
automatic substitution will eventually replace him, but a manual swap
is also available. The advisor weighs the expected gain of manually
starting an unplayed bench player against the expected loss from
forfeiting automatic substitutions for any remaining unplayed
starters in the round, and recommends the action with the higher
expectation.
